%% supplementary.tex — v2: honest circularity, direction consistency, u values, DBTL architecture

\section*{Supplementary Materials}

\subsection*{Table S1: RNACTM consistency verification}

\small Note: half-life consistency values reflect parameterization by design, not independent prediction. Rate constants are derived from the same literature sources listed below; the $<$4.1\% deviations reflect numerical precision of the ODE solver, not predictive accuracy. The most informative metric is the expression window deviation (17\%), which captures dynamics not directly parametrized. Sensitivity: $\pm$20\% rate constant variation changes half-life by $\leq$8.2\%.

\begin{tabular}{llllll}
\toprule
\textbf{Parameter} & \textbf{Literature} & \textbf{Simulated} & \textbf{Deviation} & \textbf{Source} & \textbf{Nature} \\
\midrule
RNA half-life (unmodified) & 6.0 h & 6.24 h & 4.0\% & Wesselhoeft 2018 & By design \\
RNA half-life (m6A) & 10.8 h & 11.24 h & 4.1\% & Chen 2019 & By design \\
RNA half-life ($\Psi$) & 15.0 h & 15.61 h & 4.1\% & Liu 2023 & By design \\
Endosomal escape & 1--5\% & 5.2\% & slightly above range & Gilleron 2013 & Parametrized \\
Tissue (liver) & 80\% & 80\% & 0\% & \citep{paunovska2018} & Direct assignment \\
Expression window & $\sim$48 h & 40 h & 17\% & Wesselhoeft 2018 & Emergent property \\
\bottomrule
\end{tabular}

\subsection*{Table S2: MOE baseline comparison}

\small Drug results on 91\,K samples, 6 targets. Binding R$^2{=}0.95$ (Ridge); efficacy R$^2{=}0.60$ (MOE). The high binding R$^2$ reflects that molecular binding is well-captured by physicochemical features; efficacy is a composite clinical endpoint that is inherently harder to predict.

\begin{tabular}{llllll}
\toprule
\textbf{Epitope ($N$=300, 5-fold CV)} & \textbf{R$^2$} & \textbf{MAE} & \textbf{Drug (91K, 6 targets)} & \textbf{R$^2$} & \textbf{MAE} \\
\midrule
Ridge & 0.533 & 0.639 & Ridge (binding) & 0.949 & 0.041 \\
HGB & 0.794 & 0.409 & Ridge (efficacy) & 0.542 & 0.197 \\
RF & 0.704 & 0.498 & MOE (binding) & 0.952 & 0.039 \\
GBR & 0.664 & 0.527 & MOE (efficacy) & 0.602 & 0.163 \\
MLP & 0.338 & 0.771 & & & \\
MOE & 0.819 & 0.389 & & & \\
\bottomrule
\end{tabular}

MAE improvement over Ridge: epitope 39\% ($p{<}0.001$), drug binding 4.9\%, drug efficacy 17.3\%. Removing Morgan FP (2048 dims) improves binding R$^2$ from 0.67 to 0.95.

\subsection*{Table S3: Gene signature detailed results}

\small C-index 0.52 is near-random; this dimension is flagged as exploratory in the 5D framework. Log-rank $p{<}0.001$ reflects large-sample tertile stratification, not prognostic power.

\begin{tabular}{ll}
\toprule
\textbf{Metric} & \textbf{Value} \\
\midrule
Overall C-index & 0.524 \\
C-index (TCGA-BRCA, $N$=1086) & 0.566 \\
C-index (METABRIC, $N$=1978) & 0.537 \\
Spearman $r$ (target availability vs OS) & 0.472, $p{=}9.04{\times}10^{-170}$ \\
Log-rank chi2 & 175.96, $p{<}0.001$ \\
High-target death rate & 0.568 \\
Low-target death rate & 0.202 \\
Per-gene C-index: TROP2 & 0.499 \\
Per-gene C-index: NECTIN4 & 0.550 \\
Per-gene C-index: LIV-1 & 0.518 \\
Per-gene C-index: B7-H4 & 0.519 \\
Cox regression & Failed (collinearity) \\
\bottomrule
\end{tabular}

\subsection*{Table S4: 5D weight sensitivity}

Under 1\,000 $\pm$50\% weight perturbations:
\begin{tabular}{llll}
\toprule
\textbf{Modification} & \textbf{Go rate} & \textbf{Conditional rate} & \textbf{No-Go rate} \\
\midrule
$\Psi$ & 100\% & 0\% & 0\% \\
m6A & 0\% & 99.9\% & 0.1\% \\
Unmodified & 0\% & 100\% & 0\% \\
\bottomrule
\end{tabular}

Most critical dimension: Kinetics (removing shifts $\Psi$ from Go to Conditional, $\Delta{-4.6\%}$).

\subsection*{Per-dimension uncertainty values}

\begin{tabular}{lllll}
\toprule
\textbf{Dimension} & \textbf{Base weight} & \textbf{$u_i$} & \textbf{$(1-u_i)^2$} & \textbf{Effective weight} \\
\midrule
Clinical (drug) & 0.30 & 0.15 & 0.7225 & 0.2168 \\
Binding (epitope) & 0.20 & 0.15 & 0.7225 & 0.1445 \\
Kinetics (PK) & 0.15 & 0.20 & 0.64 & 0.096 \\
Gene Signature & 0.15 & 0.50 & 0.25 & 0.0375 (exploratory) \\
circRNA: half-life/immune & 0.20 & 0.15 & 0.7225 & 0.1445 \\
circRNA: miRNA/RBP/transl & (within circRNA) & 0.80 & 0.04 & $\sim$0.006 \\
\bottomrule
\end{tabular}

Gene Signature ($u{=}0.50$) reflects near-random C-index; miRNA/RBP/translation ($u{=}0.80$) reflect no independent validation.

\subsection*{Table S5: Evolution round-by-round results}

\begin{tabular}{llllllll}
\toprule
\textbf{Round} & \textbf{Candidates} & \textbf{Top-k} & \textbf{Best reward} & \textbf{Avg toxicity} & \textbf{Avg QED} & \textbf{Pareto front} & \textbf{Gated out} \\
\midrule
1 & 48 & --- & 0.712 & 0.423 & 0.645 & 8 & 0 \\
2 & 48 & 12 & 0.823 & 0.356 & 0.712 & 14 & 6 \\
3 & 48 & 12 & 0.891 & 0.289 & 0.756 & 19 & 8 \\
4 & 48 & 12 & 0.934 & 0.234 & 0.778 & 23 & 5 \\
5 & 48 & 12 & 0.947 & 0.198 & 0.782 & 27 & 4 \\
\bottomrule
\end{tabular}

circRNA-specific operations: IRES optimization +0.067, modification switching +0.056 ($\Psi$ most beneficial), backbone mutation +0.045, BRICS crossover +0.034, UTR rearrangement +0.023.

Risk gating: adaptive quantile (q=0.75) reduces high-tox-high-efficacy candidates from 23.4\% to 3.1\%, maintains efficacy 0.702, toxicity-efficacy correlation shifts from +0.34 to $-0.18$.

\subsection*{Table S6: MHC-II allele coverage}

37 MHC-II alleles (DRB1/DQ/DP), encoded via NetMHCIIpan-4.0 34-position pseudo sequences (680 dims) + allele one-hot (37 dims) + binding position encoder (230 dims, anchors P1/P4/P6/P7/P9) = 947 total dims. Auto-detect routes HLA-A/B/C to MHC-I (1018 dims) and HLA-DR/DQ/DP to MHC-II (947 dims).

\subsection*{Version history}

Confluencia v1.0 (2026-01) provided a monolithic frontend integrating GNN molecular prediction, PINN/PDE-constrained compartment simulation, multiscale GNN-PINN cascade, cross-attention protein-ligand docking, ADMET screening, clinical trial simulation, and four-target gene signature scoring. v2.0 (2026-04) restructured into three independent modules (Drug, Epitope, circRNA) with a shared library, and added RNACTM six-compartment simulation, five innate immune pathway scoring, circRNA functional axes, uncertainty-adaptive 5D evaluation, cross-modality Go/Conditional/No-Go decisions, and circRNA-specific evolutionary operations. v2.5 (2026-05) finalized module separation with bilingual UI. v2.6 (2026-05) added allele-aware MHC training (AUC 0.74→0.80), R package (27 functions), VS Code extension (11 commands), plugin system, Confluencia Hub (offline-first federated model sharing), and open training pipelines enabling the DBTL Learn stage.

\subsection*{Immunogenicity direction consistency on real circRNA sequences}

Five modification conditions scored on three real circRNA sequences (circFOXO3, circHIPK3, circMBOAT2) plus one AU-rich synthetic. Direction consistency ($\Psi$-100\% $<$ $\Psi$-25\% $<$ unmodified) holds on all 4 sequences (4/4 pass). Per-sequence results:

\begin{tabular}{lllll}
\toprule
\textbf{Sequence} & \textbf{Unmod} & \textbf{$\Psi$-25\%} & \textbf{$\Psi$-100\%} & \textbf{m6A} \\
\midrule
circFOXO3 & 0.5649 & 0.5115 & 0.1928 & 0.5649 \\
circHIPK3 & 0.3433 & 0.3237 & 0.0766 & 0.3433 \\
circMBOAT2 & 0.5423 & 0.5423 & 0.1758 & 0.5423 \\
AU-rich & 0.1726 & 0.1726 & 0.0151 & 0.1726 \\
\bottomrule
\end{tabular}

$\Psi$-100\% consistently achieves lowest immunogenicity; $\Psi$-25\% achieves intermediate; unmodified highest. This constitutes direction consistency, not quantitative validation.

\subsection*{Real-data joint evaluation}

Joint evaluation on real IEDB/ChEMBL inputs (aspirin SMILES, SLYNTVATL epitope, HLA-A*02:01, 200mg BID) produces a Go recommendation but only the binding dimension returns a valid score (1.0, strong binder); clinical (0.0), kinetics (0.0, NaN in PK summary), gene signature (None), and circRNA (None) sub-dimensions fail due to pipeline errors. This demonstrates that single-sample joint evaluation works but 4/5 dimensions require further integration work---a known limitation that the adaptive framework handles via $(1{-}u)^2$ downweighting.

\subsection*{ESM-2 overfitting experiment}

On $N{=}300$ training samples, ESM-2 (650M parameter protein language model; Lin et al. 2023) produced severe overfitting: training R$^2{>}0.95$ but validation R$^2{<}0.3$ across all 5 folds, with MAE increasing 2--3$\times$. This motivated the sample-size-adaptive MOE design using well-regularized sklearn models.

\subsection*{Immunogenicity direction consistency}

Five modification conditions scored by Confluencia's heuristic pathway weights agree in direction with literature IFN rankings ($\rho{=}0.93$ over 5 conditions, $N$ too small for statistical significance): $\Psi$ lowest, unmodified highest. This constitutes direction consistency, not quantitative validation. Pathway weights (RIG-I 35/40/20/5\%, TLR7 45/30/20/5\%, PKR 50/25/20/5\%) are author-informed heuristics derived from qualitative literature observations, not empirically calibrated.